\section{Background}
\label{sec:background}

The design of accountable, autonomous multi-agent systems requires a precise answer to a foundational question: when one agent delegates a task to another, through what structural mechanism is the chain of accountability preserved? This question becomes acute in systems that support \emph{recursive spawning} --- the capacity for an agent to decompose its assigned task by provisioning child agents that operate with delegated authority. As agentic runtimes move from single-agent to hierarchical multi-agent architectures \cite{wooldridge1995,stone2000}, the problem of managing the lifecycle of spawned agents, and of maintaining an unbroken accountability chain from every active agent back to an originating human decision, transitions from an architectural detail to a core safety requirement.

This section situates Recursive Spawning within five distinct research and engineering traditions: agent lifecycle management in multi-agent systems; accountability and provenance in distributed autonomous systems; fixed versus mutable delegation chains; hierarchical task decomposition in classical and modern AI; and the BX3 Framework's three-layer model as an architectural substrate for immutable parent chains. The section establishes the conceptual foundations upon which the Recursive Spawning Protocol (RSP), defined in Section~\ref{sec:rsp}, is constructed.

% -------------------------------------------------------
\subsection{Agent Lifecycle Management in Multi-Agent Systems}
% -------------------------------------------------------

Multi-agent systems (MAS) research has long recognized that the lifecycle of an agent --- its creation, activation, operation, suspension, and termination --- is a first-class design concern, not merely an implementation detail. The taxonomy introduced by Stone and Veloso \cite{stone2000} classifies multi-agent systems along axes of agent heterogeneity (homogeneous versus heterogeneous populations), agent knowledge (perfect versus imperfect information), agent access to the environment (fully versus partially accessible), and the nature of agent interactions (cooperative versus competitive). Within this taxonomy, recursive spawning is most relevant to cooperative, homogeneous systems operating in partially accessible environments --- precisely the regime in which modern agentic runtimes such as AutoGPT \cite{autogpt2023}, LangChain \cite{langchain2023}, and Microsoft's AutoGen \cite{autogen2023} operate.

The management of agent lifecycles in such systems must address at minimum four distinct concerns: \emph{creation} (how an agent is instantiated and initialized with authority), \emph{delegation} (how authority is transmitted from parent to child), \emph{surveillance} (how the parent monitors the child's execution state), and \emph{termination} (how the child exits and returns results or an error state). Ferber and M\"uller \cite{ferber1996} formalized these as interactions between an agent's lifecycle state machine and the organizational context in which it operates, arguing that organizational roles --- not just individual agent states --- must be modeled to ensure system-wide coherence. This insight is directly relevant to Recursive Spawning: the spawn chain is not merely a tree of process instances; it is an organizational hierarchy of delegated authority, and each node's lifecycle must be governed with respect to both its local state and its position in that hierarchy.

The concept of agent delegation was formalized by Yim et al. \cite{yim1997} as a structured transfer of goals, plans, and execution authority from one agent to another, distinguished from simple message passing by the explicit surrender of control. Subsequent work by Padgham and Thiel \cite{padgham2001} demonstrated that incomplete lifecycle management --- particularly the failure to distinguish between suspension and termination states --- introduces subtle but systemic failures in agent systems, where orphaned or zombie agents continue to consume resources and emit effects long after their parent has exited. These findings establish that lifecycle management is not a single design decision but an ongoing architectural discipline: the system must maintain coherent authority boundaries at every state transition throughout the entire spawn chain.

More recent work on multi-agent orchestration has surfaced the specific problem of \emph{unbounded spawning}. The emergence of large language model-based agents capable of tool use and self-prompted task decomposition has made the uncontrolled proliferation of child agents a practical concern rather than a theoretical one. Zhuge et al. \cite{zhuge2024} document the ``agent chain problem'': LLM-based agents that recursively spawn children to handle sub-tasks, without a defined termination condition or depth ceiling, can produce chains of dozens of intermediate agents that collectively exceed human capacity for oversight and whose aggregate authority may differ substantially from what any single agent was authorized to exercise. This problem is not addressed by classical lifecycle management frameworks, which assume well-defined, static agent roles rather than dynamically provisioned, task-specific child agents. Recursive Spawning is designed specifically to address the dynamic spawning regime.

% -------------------------------------------------------
\subsection{Accountability and Provenance in Distributed Systems}
% -------------------------------------------------------

Accountability in autonomous systems has two structurally distinct dimensions: \emph{attribution} (after an event, who or what can be identified as the cause of the outcome) and \emph{enforcement} (before or during an event, what mechanism ensures the system behaves in an accountable way) \cite{bansal2019}. Attribution is a forensic concern; enforcement is an architectural concern. Most existing accountability frameworks address attribution retrospectively --- through logging, audit trails, and post-hoc analysis --- without constraining the system's runtime behavior to maintain the conditions under which those audit trails remain meaningful \cite{ristenpart2009}.

This distinction is particularly consequential in distributed systems where agents operate asynchronously, across process boundaries, and with partial visibility into each other's state. The problem is compounded when agents can spawn children, because each spawn event creates a new locus of authority that may not be traceable to the original authorization. An agent $A$ is authorized to act on behalf of a human principal $H$; $A$ spawns child $B$; $B$ spawns child $C$; $C$ makes a consequential decision. Who is accountable for $C$'s decision? In conventional architectures, the answer depends on policy documents and implicit trust relationships that were established at design time and may not accurately reflect the actual chain of delegation that emerged at runtime.

Provenance tracking in distributed systems addresses the attribution problem directly. The Open Provenance Model (OPM) \cite{moreau2011} established a formal vocabulary for describing the causal relationships between entities in a distributed computation: artifacts, processes, and agents. OPM's key insight is that provenance must be captured \emph{lifecycle} --- at the time of each relationship creation, not reconstructed later from logs --- and that provenance records must be immutable once written. The Linked Data Platform provenance model extends this to web-scale systems, emphasizing that provenance chains can span organizational boundaries and must therefore be expressed in interoperable, machine-readable formats \cite{nyberg2012}.

Immutable provenance is a necessary but not sufficient condition for runtime accountability. A provenance record that is written immutably but not consulted at decision time cannot enforce anything; it can only document after the fact. The contribution of Recursive Spawning is to make provenance records \emph{operative} rather than merely documentary: the immutable parent pointer stored in the spawn warrant is not only a forensic artifact but a runtime enforcement mechanism. A child agent cannot act without a verified warrant, and its warrant cannot be verified without traversing the full chain back to a human anchor. The chain is the enforcement mechanism, not just the record of it.

% -------------------------------------------------------
\subsection{Fixed Versus Mutable Delegation Chains}
% -------------------------------------------------------

Delegation chains in multi-agent systems take two structurally different forms, with profoundly different accountability properties. In a \emph{fixed delegation chain}, the parent-child relationship is established at spawn time and is immutable thereafter: the child's parent pointer cannot be changed, the child's authority cannot be re-delegated except through the defined spawn protocol, and the chain terminates at a predetermined anchor (typically a human principal). In a \emph{mutable delegation chain}, the parent-child relationship is treated as a runtime state that can be updated, redirected, or reassigned as conditions change.

Mutable delegation chains offer greater flexibility: a child agent can be reparented to a different parent if the organizational structure changes, if a parent fails, or if load balancing requires reassignment. This flexibility, however, comes at a significant accountability cost. If the parent-child relationship is mutable at runtime, then any single action by a child agent may have been authorized by one parent at one moment and a different parent at another --- and the effective authorization for any given action may not be the authorization that was recorded at spawn time. The chain of provenance becomes conditional on the state of a mutable mapping, rather than fixed by an immutable record. This introduces what we term the \emph{delegatee ambiguity problem}: after an event, it may be impossible to determine which principal effectively authorized the action, because the delegation chain has been rewritten since the action occurred.

Fixed delegation chains sacrifice this flexibility for a stronger accountability guarantee. Each child's authority is traceable to a single, immutable parent pointer, and that pointer chains unbroken back to a human anchor. There is no ambiguity about effective authorization: the warrant that activated the child encodes the complete delegation chain as it existed at the moment of activation, and that chain cannot be altered retroactively. The cost is that reparenting requires a formal re-spawn: the child must be terminated, a new spawn event must be initiated, and a new warrant must be issued. This is more expensive than a pointer update but produces an unambiguous, immutable authorization record for the new chain.

The BX3 Framework's Fact Layer enforces fixed delegation chains as a matter of architectural principle. The immutable parent pointer is not a configurable policy option; it is a structural requirement of the Fact Layer's authorization model. We argue that in accountable autonomous systems, the accountability benefit of fixed chains substantially outweighs the flexibility cost, for two reasons. First, the organizational cost of re-spawning a child agent is typically small compared to the organizational cost of an ambiguous accountability chain after a consequential failure. Second, genuine organizational restructuring that requires reparenting is rare in practice; most mutable reparenting requests in conventional systems are actually attempts to route around a spawn protocol that was established to prevent precisely the kind of dynamic authority expansion that mutable chains permit.

% -------------------------------------------------------
\subsection{Hierarchical Task Decomposition in AI}
% -------------------------------------------------------

Hierarchical task decomposition is the problem of breaking a high-level goal or directive into a structured sequence of sub-goals that can be executed by specialized agents or modules. It is one of the foundational techniques in classical AI planning and has experienced a significant resurgence with the advent of large language model-based agents.

In classical AI, hierarchical task network (HTN) planning \cite{erol1994,ghallab2004} formalizes decomposition as a recursive process: a task is decomposed into subtasks by applying methods from a predefined library; each subtask is itself either primitive (directly executable) or compound (requiring further decomposition). HTN planning's key theoretical property is that it reduces the search space of plan generation by exploiting domain knowledge encoded in the method library: instead of searching over all possible action sequences, the planner searches over decomposition trees, which are typically much smaller. The completeness of HTN planning depends on the completeness of the method library; if no method applies to a compound task, decomposition fails and the planner must backtrack.

Modern LLM-based agents perform a form of informal, natural-language HTN decomposition: given a high-level directive, the agent generates a decomposition into subtasks, assigns them to tool calls or child agents, and synthesizes results. Chain-of-thought prompting \cite{wei2022} and its descendants (tree-of-thought \cite{yao2023}, graph-of-thought \cite{besta2024}) formalize this process as an explicit reasoning trace, making decomposition decisions visible and auditable rather than opaque. The ReAct paradigm \cite{yao2022} extends this by interleaving reasoning traces with environment observations, allowing the agent to adapt its decomposition dynamically as task conditions change.

What connects classical HTN and modern LLM-based decomposition is the spawn tree: a hierarchical structure in which each node represents a subtask, and edges represent the decomposition relationship. In both traditions, the tree is typically mutable and implicit --- the decomposition decisions are not recorded as first-class authorization artifacts, and child nodes are not distinguished from parent nodes in terms of their accountability status. A classical HTN planner's subtasks are subroutine calls; they do not have independent accountability chains. An LLM agent's spawned sub-agents are independent processes, but the chain of authorization is implicit in the agent's context window rather than explicit in an immutable warrant structure.

Recursive Spawning adopts the hierarchical decomposition paradigm from both traditions but adds the critical accountability layer that is absent from both: each node in the decomposition tree is a first-class accountability entity with an immutable parent pointer, a defined termination condition, and a cryptographically verifiable chain back to a human anchor. Decomposition is not merely a planning technique; it is an authorization event that creates a new accountable entity. The decomposition tree is therefore not just a planning artifact but an organizational chart of delegated authority.

% -------------------------------------------------------
\subsection{The BX3 Framework as Substrate for Immutable Parent Chains}
% -------------------------------------------------------

The BX3 Framework provides the architectural substrate within which Recursive Spawning operates. As described in \cite{bx3framework2026}, the BX3 Framework organizes any autonomous node into three immutable functional layers --- the \textit{Purpose Layer}, the \textit{Bounds Engine}, and the \textit{Fact Layer} --- each defined by the properties it must maintain rather than by the type of actor occupying it. The Purpose Layer carries intent, judgment, and final human accountability; the Bounds Engine carries constrained execution within a defined operating range; the Fact Layer carries deterministic enforcement and forensic auditability. The upstream accountability guarantee of the framework states that when any node encounters a state it cannot resolve within its bounds, accountability escalates recursively upward through the system hierarchy, bypassing all machine actors, until it reaches a human anchor \cite{bx3framework2026}.

The Fact Layer plays the central role in Recursive Spawning. It is the Fact Layer that creates the immutable parent pointer, issues the spawn warrant, and maintains the forensic ledger that records every spawn event. The Fact Layer's determinism is what makes the parent pointer immutable: once written, the pointer cannot be updated because the Fact Layer does not permit modifications to its authorization records. This immutability is not a policy choice; it is an architectural property of the deterministic enforcement layer.

The three-layer model provides the structural basis for the upstream accountability guarantee in the spawn context specifically. When a child agent is activated via a spawn warrant, its authority is bounded by the operating range defined by the parent node's Bounds Engine. If the child encounters a condition outside its provisioned scope, it triggers the Bailout Protocol \cite{bailoutprotocol2026} --- the upward exception propagation mechanism of the BX3 Framework --- which bypasses all intermediate machine actors and routes the exception to the nearest human accountability anchor. The spawn chain is therefore not only an authorization structure but an escalation structure: the same chain that carries delegated authority upward carries exceptions downward, and when the Bounds Engine cannot resolve them, carries accountability upward to the human who authorized the chain in the first place.

This coupling of authorization and escalation is the key architectural insight of Recursive Spawning. In conventional systems, authorization and escalation are separate concerns: a spawning decision is made by the parent agent, and an escalation decision is made independently when an exception occurs. In the BX3 Framework, both decisions are governed by the same immutable chain structure. The chain that authorized the child is the chain along which exceptions propagate when the child exceeds its bounds. There is no separate escalation infrastructure; the spawn chain is the escalation chain, by design.

The result is an accountability architecture in which the Question --- \emph{who is responsible for this agent's actions?} --- is always answerable by traversing the immutable parent chain back to a human anchor, and in which every traversal of that chain is a structurally enforced property of the runtime rather than a forensic reconstruction performed after the fact. The Recursive Spawning Protocol, defined formally in Section~\ref{sec:rsp}, is the mechanism by which this property is operationalized.
